Los puntos de la figura \ref{fig:grafica} caen casi perfectamente sobre una línea recta. Eso es lo importante, porque significa que el cociente entre los senos se mantiene constante sin importar el ángulo de entrada, y esa es precisamente la afirmación de la ley de Snell.

¿Y el valor del índice?

Ahí la historia cambia un poco. La pendiente dio 1.396 y el valor de referencia para el acrílico es 1.49, de modo que el error porcentual es
\begin{equation}
E = \frac{\left| 1.49 - 1.396 \right|}{1.49} \times 100 = 6.3 \text{ por ciento},
\label{eq:error}
\end{equation}
dentro del rango de 5 a 10 por ciento con el que se construyeron los datos. Nótese que todos los puntos quedan por debajo de la línea punteada de referencia, y no repartidos a ambos lados: el error no es solo ruido al azar, sino un corrimiento parejo de los ángulos de refracción hacia valores un poco más grandes.

En un experimento real este tipo de corrimiento tendría causas bastante concretas. Bastaría con que la cara plana de la lente no quedara justo sobre la línea COMPONENT, o con que su centro estuviera unos milímetros fuera del centro de la mesa, para que el rayo ya no saliera por la cara curva a lo largo de un radio y se doblara una segunda vez. A eso se suman el grosor del rayo, que en la escala ocupa casi un grado, y la dificultad de leer en la oscuridad.

Otro detalle. Con el índice obtenido, el ángulo crítico de la ecuación \eqref{eq:critico} sería de 45.8 grados en vez de los 42.2 grados que corresponden a $n = 1.49$; un error de unas cuantas centésimas en el índice se nota bastante en este ángulo.

En conclusión, la relación $\sin\theta_i = n \sin\theta_t$ describe muy bien el comportamiento del rayo, con un ajuste lineal de $R^2 = 0.9986$, y de la pendiente se obtiene un índice de refracción de 1.396 para el acrílico, 6.3 por ciento por debajo del valor de referencia. Hay que insistir en que estos números vienen de datos simulados, por lo que confirman cómo se analiza la práctica, pero no sustituyen una medición hecha en el laboratorio.

Queda pendiente repetirla con el equipo real.
